\section{Introduction}
\label{sec:introduction}

A judge orders the government to buy a cancer drug within \BPurgentPlanLow--\BPurgentPlanHigh{} days. The procurement officer has no time to aggregate demand across other purchases, no time to research market prices, and faces personal fines if delivery fails. The drug arrives, but the cost margin between this transaction and the same agency's ordinary purchase of the same drug runs from a few percent at the typical urgent purchase to as much as a third more for purchases that are court-mandated rather than initiated through an internal administrative channel.

This paper provides the first quantitative decomposition of that cost margin. Courts and regulatory agencies are alternative instruments of social control \citep{glaeser2003rise}, but court intervention in procurement compresses planning timelines and distorts incentives in ways that raise prices and reduce competition \citep{coviello2018court}. In health care, the stakes are high: judicial mandates directly determine how---and at what price---governments acquire essential medicines \citep{wang2015right, ferraz2009right}. Brazil's 1988 Constitution enshrines health as a ``right of all and a duty of the state,'' and courts have granted virtually every medication request---\BPlitClaimsBR{} first-instance claims nationally in \BPcnjClaimsYear, a \BPspGrowthSince{} increase in S\~ao Paulo since \BPspGrowthBaseYear{} \citep{cnj2019judicializacao}.

We use bid-level data covering all pharmaceutical procurement by the S\~ao Paulo State Department of Health (SES/SP) from \BPsampleStartYear{} through \BPsampleEndYear: over \BPnPOIfull{} purchase-offer-item observations conducted through the state's electronic procurement platform (BEC). Court orders and administrative requests partition purchases into \textit{ordinary} (standard timeline) and \textit{urgent} (compressed timeline). Within urgent purchases, a further institutional distinction proves useful: \textit{administrative} requests, instituted by SES/SP in \BPsampleStartYear{}, carry no penalty for procurement failure, while \textit{litigated} purchases expose officials to fines, personal liability, and fund seizure. Comparing these two types---which share planning constraints but differ in penalty exposure---is what we call the ``under the gun'' (UTG) comparison. \emph{Our identification is descriptive throughout. We decompose a within-item gap between two procurement regimes that share planning constraints but differ in penalty exposure; we do not estimate a counterfactual price under hypothetical sanction removal. That object is unidentified in our setting because the administrative subsample is itself selected by a scientific committee on cost-effectiveness criteria.} What we provide is a quantitative decomposition of where the cost margin between regimes comes from.

Across the full sample, urgent purchases carry negotiated prices \BPnegPctHeadline{} higher, reference prices \BPrefPctPreferred{} higher, and attract \BPfirmsPctHeadline{} fewer bidders---yet succeed \textit{more often} (\BPsuccessPP), consistent with officials accepting worse terms to avoid failure. The UTG comparison is sharper: litigated purchases cost \BPutgPctRange{} more than administrative ones in the same item under the same fixed effects. We then decompose this UTG umbrella into three distinct channels. \textit{Demand fragmentation} (C1) does most of the work: orders shrink under judicial pressure, and the bulk-discount channel mechanically accounts for essentially all of the UTG gap and roughly half of the urgency premium across the full sample. \textit{Reduced participation} (C2) explains the residual urgency premium across the full sample: even at identical order size, urgent tenders attract about \BPfirmsPanelBpct{} fewer bidders, leaving thinner markets and worse negotiated prices. \textit{Supplier composition shift} (C3) explains the rest of the across-sample premium: the supplier mix differs between urgent and ordinary regimes; within firm-buyer-item triples observed in both regimes, the within-firm price under urgency is, if anything, slightly lower, ruling out a within-firm markup story.

Three contributions follow. First, we provide the first quantitative decomposition of the procurement cost of judicial enforcement, contributing to a large literature on government spending efficiency and procurement design \citep{bandiera2009active, best2023procurement, decarolis2018buyer, lewisfaupel2016can, szucs2024discretion}. The decomposition relies on a methodological discipline that explicitly avoids post-treatment-bias: we read controls on quantity and supplier identity as channel-decomposing rather than as direct-effect estimators \citep{acharya2016explaining}. Second, the UTG comparison---exploiting within-urgent variation in penalty exposure---offers direct evidence that accountability mechanisms can distort bureaucratic decisions, echoing a long-standing theoretical prediction \citep{prendergast2007bureaucrats} and complementing recent field evidence on bureaucratic effort, autonomy, and management \citep{rasul2018management, williams2021bureaucratic, bandiera2021allocation}. Third, in contrast to \citet{coviello2018court}, where slow courts \emph{extend} procurement timelines, we study injunctions that \emph{compress} them, adding to the evidence on ex ante design in auction performance \citep{coviello2018bidding, baltrunaite2021discretion} and on competition design under contractual incompleteness \citep{carril2026competition}. We complement the legal and public-health literatures on the judicialization of health in Latin America \citep{wang2015right, ferraz2009right, biehl2009will}, which have focused on legal and clinical dimensions with limited attention to procurement costs.

The remainder of the paper proceeds as follows. Section~\ref{sec:institutional} describes the institutional setting. Section~\ref{sec:data} presents the data. Section~\ref{sec:empirical} details the empirical strategy. Section~\ref{sec:results} presents results. Section~\ref{sec:conclusion} concludes.
